\section{Implications for clinical-simulation validation}
\label{sec:validation}

The audit measures a failure of the recorded reporting claim; the synthetic
illustration explains why target-population alignment is a separate requirement.
Three established checks make these distinctions operational.

\paragraph{State the target before constructing its reference.}
Distinguish a population treatment effect, a latent finite-cohort effect and
an observed-data contrast. Name the population and the averaging weights.
A reference need not be independent of the simulated dataset: $\tau_E$ is
necessarily a property of the enrolled records. The requirement is that it
represents the intended target, rather than merely reproducing the estimator.

\paragraph{Evaluate the interval and reporting policy jointly.}
Assess the whole decision rule using the interval that would actually
be reported. Record bias and error against the chosen target, null rejection,
detection at the specified effect, and abstention frequency. Report both
conditional interval performance among returned answers and how often a valid
answer is returned. A successful effect-detection rate can accompany excessive
null rejection, so the former alone cannot validate an interval. For a nominal 95\% interval, null rejection and null noncoverage are
the same quantity; alternative coverage and detection answer different questions.

\paragraph{Separate rule selection from assessment and transfer.}
Separate selecting a rule from assessing it. If a scalar threshold is
intended to guarantee performance above that threshold, qualification must
address the evaluated upper range rather than only its first passing point.
Reserve independent simulation draws for the final assessment, report rate
uncertainty and the tested range, and distinguish this from transfer to new
patients or cohorts. Increasing resamples from an empirical pool improves
conditional numerical precision; it does not broaden that pool's biological
support. For an external-control application, the additional question is
comparability across data sources, which requires direct study-level validation
such as the held-out comparisons of \citet{ventz2019}.
